% ЗГЕНЕРОВАНО scripts/build_format_v2_chapter.py — не редагувати вручну.
% Джерело: books/ashtavakra-gita/poetychnyi/rozdil-16.md

\poemverse{16.1}{Розповідай чи слухай, любий, \\ хоч би скільки різних наук, — \\ спокою тобі не буде, \\ поки не забудеш усе.}
\poemverse{16.2}{Насолоджуйся, дій або зосереджуйся, \\ обізнаний, — \\ та по-справжньому смакуватиме тобі \\ лише тоді, \\ коли розум відкине всі надії.}
\poemverse{16.3}{Від натуги всі нещасні, \\ і ніхто з них цього не бачить. \\ Блаженний же досягає вгамування \\ самим лише цим настановленням.}
\poemverse{16.4}{Кому важко навіть повіки звести — \\ той перший серед лінивих: \\ немає щастя ні йому, \\ ні комусь іншому такому.}
\poemverse{16.5}{Коли розум звільниться від пари \\ це я зробив — а це ні, \\ він стане байдужим \\ і до дгарми, і до артги, \\ і до ками, і до мокші.}
\poemverse{16.6}{Відвернений — це той, хто ненавидить чуттєве; \\ пристрасний — той, хто до нього жадібний. \\ А хто не хапає і не відкидає — \\ той уже ні відвернений, ні пристрасний.}
\poemverse{16.7}{Розрізнення відкинути це, взяти те — \\ лише пагін на дереві сансари; \\ він росте, доки жива жадоба, \\ а жадоба — осідок життя без роздуму.}
\poemverse{16.8}{Коли дієш — народжується потяг, \\ коли спиняєшся — відраза. \\ Мудрий, вільний від цієї пари, \\ стоїть так, як стоїть дитина.}
\poemverse{16.9}{Пристрасний хоче втекти з сансари, \\ бо хоче позбутися страждання. \\ А в кого пристрасті вже немає, \\ той не страждає — \\ і навіть у сансарі не мучиться.}
\poemverse{16.10}{Якщо людина претендує навіть на звільнення \\ і водночас відчуває тіло як своє, \\ вона не знавець і не йог — \\ вона просто має свою частку страждання.}
\poemverse{16.11}{Навіть якби твоїм наставником був сам Шіва, \\ Вішну чи Брахма, — \\ спокою тобі не буде, \\ поки не забудеш усе.}
